\begin{abstract}
Quantization stores the weights of a large language model (LLM) in 8 or 4 bits instead of 16, and
formats are usually compared by their bit-width. We show that the bit-width alone predicts neither the
memory a format saves, nor its speed, nor its cost in quality. On one NVIDIA RTX 5090, we measure
eight configurations on two inference stacks (Hugging Face Transformers with bitsandbytes, and vLLM)
for five models from 3B to 32B parameters, with \NumMeasurements{} timed runs, real chat traffic and
three quality benchmarks. (1)~4-bit weights make a
7B model \WtNFSeven--\WtAWQSeven{} smaller, not $4\times$, and the total saving falls to \MemBigNF{}
as the KV cache grows. (2)~The kernel matters more than the bit-width: the same 8-bit weights decode
\GPTQeightspeedup{} faster than BF16 with vLLM's Marlin kernel but \INTeightslowdown{} slower with
bitsandbytes. (3)~Weight-only formats lose their speedup once batching makes decoding compute-bound,
as a roofline model predicts, while FP8 stays faster at every batch size. (4)~8-bit formats keep
quality, and 4-bit formats cost \MFourAbsMin--\MFourAbsMax{} MMLU points, more on smaller models.
Still, a larger 4-bit model beats a smaller 16-bit one: Qwen3-32B in 4 bits scores \BigMMLUawq\% on
MMLU against \SmallMMLUbf\% for Qwen3-8B in BF16, with only \BigMemRatio{} the weight memory. Kernel
benchmarks, a roofline model and GPU profiles explain each result.
\end{abstract}
